% !TeX root = part2.tex

\documentclass[../final.tex]{subfiles}

\begin{document}

% При отдельной сборке продолжаем нумерацию полного конспекта.
\ifSubfilesClassLoaded{\setcounter{chapter}{1}\setcounter{section}{7}}{}


% ============================================================
\section{Ферми-газ при конечных температурах}
% ============================================================

Рассмотрим сильно вырожденный ферми-газ в области
%
\begin{equation*}
    T
    \ll
    \varepsilon_F.
\end{equation*}

При конечной, но малой температуре резкая ступенька распределения
Ферми--Дирака размывается лишь в узкой области энергий порядка $T$
вблизи химического потенциала. Поэтому температурные поправки
к термодинамическим величинам удобно находить разложением
интегралов с распределением Ферми--Дирака.


% ------------------------------------------------------------
\subsection{Разложение интеграла Ферми}
% ------------------------------------------------------------

Рассмотрим интеграл
%
\begin{equation*}
    I
    =
    \int_0^\infty
    \frac{
        f(\varepsilon)\,d\varepsilon
    }{
        e^{(\varepsilon-\mu)/T}+1
    }.
\end{equation*}

Введём новую переменную
%
\begin{equation*}
    \varepsilon-\mu
    =
    Tz.
\end{equation*}

Тогда
%
\begin{equation*}
    \varepsilon
    =
    \mu+Tz,
    \qquad
    d\varepsilon
    =
    T\,dz,
\end{equation*}
%
а пределы интегрирования переходят в
%
\begin{equation*}
    \varepsilon=0
    \quad\Longrightarrow\quad
    z=-\frac{\mu}{T},
    \qquad
    \varepsilon\rightarrow\infty
    \quad\Longrightarrow\quad
    z\rightarrow\infty.
\end{equation*}

Поэтому
%
\begin{equation*}
    I
    =
    T
    \int_{-\mu/T}^{+\infty}
    \frac{
        f(\mu+Tz)
    }{
        e^z+1
    }
    \,dz.
\end{equation*}

Разобьём область интегрирования в точке $z=0$:
%
\begin{equation*}
    I
    =
    T
    \int_{-\mu/T}^{0}
    \frac{
        f(\mu+Tz)
    }{
        e^z+1
    }
    \,dz
    +
    T
    \int_0^\infty
    \frac{
        f(\mu+Tz)
    }{
        e^z+1
    }
    \,dz.
\end{equation*}

В первом интеграле сделаем замену
%
\begin{equation*}
    z
    \rightarrow
    -z.
\end{equation*}

Получаем
%
\begin{equation*}
    I
    =
    T
    \int_0^{\mu/T}
    \frac{
        f(\mu-Tz)
    }{
        e^{-z}+1
    }
    \,dz
    +
    T
    \int_0^\infty
    \frac{
        f(\mu+Tz)
    }{
        e^z+1
    }
    \,dz.
\end{equation*}

Используем тождество
%
\begin{align*}
    \frac{1}{e^{-z}+1}
    &=
    \frac{e^z}{e^z+1}
    \\[4pt]
    &=
    1
    -
    \frac{1}{e^z+1}.
\end{align*}

Следовательно,
%
\begin{align*}
    I
    &=
    T
    \int_0^{\mu/T}
    f(\mu-Tz)\,dz
    -
    T
    \int_0^{\mu/T}
    \frac{
        f(\mu-Tz)
    }{
        e^z+1
    }
    \,dz
    \\[4pt]
    &\hspace{1.2cm}
    +
    T
    \int_0^\infty
    \frac{
        f(\mu+Tz)
    }{
        e^z+1
    }
    \,dz.
\end{align*}

В первом интеграле возвращаемся к переменной энергии:
%
\begin{equation*}
    T
    \int_0^{\mu/T}
    f(\mu-Tz)\,dz
    =
    \int_0^\mu
    f(\varepsilon)\,d\varepsilon.
\end{equation*}

При сильном вырождении
%
\begin{equation*}
    \frac{\mu}{T}
    \gg
    1.
\end{equation*}

В интеграле, содержащем множитель
%
\begin{equation*}
    \frac{1}{e^z+1},
\end{equation*}
%
область больших $z$ экспоненциально подавлена. Поэтому верхний
предел $\mu/T$ можно заменить на бесконечность:
%
\begin{equation*}
    I
    =
    \int_0^\mu
    f(\varepsilon)\,d\varepsilon
    +
    T
    \int_0^\infty
    \frac{
        f(\mu+Tz)-f(\mu-Tz)
    }{
        e^z+1
    }
    \,dz.
\end{equation*}

Разложим функции в ряд около $\varepsilon=\mu$:
%
\begin{align*}
    f(\mu+Tz)
    &=
    f(\mu)
    +
    Tz\,f'(\mu)
    +
    \frac{(Tz)^2}{2}f''(\mu)
    +
    \frac{(Tz)^3}{6}f'''(\mu)
    +
    \ldots,
    \\[6pt]
    f(\mu-Tz)
    &=
    f(\mu)
    -
    Tz\,f'(\mu)
    +
    \frac{(Tz)^2}{2}f''(\mu)
    -
    \frac{(Tz)^3}{6}f'''(\mu)
    +
    \ldots.
\end{align*}

При вычитании чётные степени сокращаются:
%
\begin{equation*}
    f(\mu+Tz)-f(\mu-Tz)
    =
    2Tz\,f'(\mu)
    +
    \frac{(Tz)^3}{3}f'''(\mu)
    +
    \ldots.
\end{equation*}

Ограничиваясь поправкой порядка $T^2$, получаем
%
\begin{equation*}
    I
    =
    \int_0^\mu
    f(\varepsilon)\,d\varepsilon
    +
    T^2f'(\mu)
    \int_0^\infty
    \frac{
        2z\,dz
    }{
        e^z+1
    }.
\end{equation*}

Известный интеграл равен
%
\begin{equation*}
    \int_0^\infty
    \frac{
        z\,dz
    }{
        e^z+1
    }
    =
    \frac{\pi^2}{12},
\end{equation*}
%
поэтому
%
\begin{equation*}
    \int_0^\infty
    \frac{
        2z\,dz
    }{
        e^z+1
    }
    =
    \frac{\pi^2}{6}.
\end{equation*}

Таким образом,
%
\[
    \rboxed{
        \int_0^\infty
        \frac{
            f(\varepsilon)\,d\varepsilon
        }{
            e^{(\varepsilon-\mu)/T}+1
        }
        =
        \int_0^\mu
        f(\varepsilon)\,d\varepsilon
        +
        \frac{(\pi T)^2}{6}
        f'(\mu)
    }.
\]

Это разложение применимо при
%
\begin{equation*}
    T
    \ll
    \mu
    \sim
    \varepsilon_F.
\end{equation*}


% ------------------------------------------------------------
\subsection{\texorpdfstring{$\Omega$}{Omega}-потенциал}
% ------------------------------------------------------------

Для нерелятивистского идеального ферми-газа
%
\begin{equation*}
    \Omega
    =
    -\frac{2}{3}AV
    \int_0^\infty
    \frac{
        \varepsilon^{3/2}\,d\varepsilon
    }{
        e^{(\varepsilon-\mu)/T}+1
    }.
\end{equation*}

Воспользуемся полученной формулой, положив
%
\begin{equation*}
    f(\varepsilon)
    =
    \varepsilon^{3/2}.
\end{equation*}

Тогда
%
\begin{equation*}
    f'(\varepsilon)
    =
    \frac{3}{2}
    \sqrt{\varepsilon}.
\end{equation*}

Следовательно,
%
\begin{align*}
    \int_0^\infty
    \frac{
        \varepsilon^{3/2}\,d\varepsilon
    }{
        e^{(\varepsilon-\mu)/T}+1
    }
    &=
    \int_0^\mu
    \varepsilon^{3/2}\,d\varepsilon
    +
    \frac{(\pi T)^2}{6}
    \frac{3}{2}
    \sqrt{\mu}
    \\[4pt]
    &=
    \frac{2}{5}
    \mu^{5/2}
    +
    \frac{(\pi T)^2}{4}
    \sqrt{\mu}.
\end{align*}

Поэтому
%
\begin{align*}
    \Omega
    &=
    -\frac{2}{3}AV
    \left[
        \frac{2}{5}\mu^{5/2}
        +
        \frac{(\pi T)^2}{4}
        \sqrt{\mu}
    \right]
    \\[6pt]
    &=
    -\frac{2}{3}AV
    \frac{2}{5}\mu^{5/2}
    \left[
        1
        +
        \frac{5}{8}
        \left(
            \frac{\pi T}{\mu}
        \right)^2
    \right].
\end{align*}

Таким образом,
%
\[
    \rboxed{
        \Omega
        =
        -\frac{2}{3}AV
        \frac{2}{5}\mu^{5/2}
        \left[
            1
            +
            \frac{5}{8}
            \left(
                \frac{\pi T}{\mu}
            \right)^2
        \right]
    }.
\]

Термодинамический дифференциал $\Omega$ имеет вид
%
\begin{equation*}
    d\Omega
    =
    -S\,dT
    -
    p\,dV
    -
    N\,d\mu.
\end{equation*}

Следовательно,
%
\begin{equation*}
    N
    =
    -
    \left(
        \frac{\partial\Omega}{\partial\mu}
    \right)_{T,V}.
\end{equation*}

Дифференцируя полученное выражение, находим
%
\begin{align*}
    N
    &=
    \frac{2}{3}AV
    \left[
        \mu^{3/2}
        +
        \frac{(\pi T)^2}{8\sqrt{\mu}}
    \right]
    \\[4pt]
    &=
    \frac{2}{3}AV
    \mu^{3/2}
    \left[
        1
        +
        \frac{1}{8}
        \left(
            \frac{\pi T}{\mu}
        \right)^2
    \right].
\end{align*}

Итак,
%
\[
    \rboxed{
        N
        =
        \frac{2}{3}AV
        \mu^{3/2}
        \left[
            1
            +
            \frac{1}{8}
            \left(
                \frac{\pi T}{\mu}
            \right)^2
        \right]
    }.
\]


% ------------------------------------------------------------
\subsection{Химический потенциал}
% ------------------------------------------------------------

При $T=0$
%
\begin{equation*}
    \mu
    =
    \varepsilon_F,
\end{equation*}
%
поэтому
%
\begin{equation*}
    N
    =
    \frac{2}{3}AV
    \varepsilon_F^{3/2}.
\end{equation*}

Поскольку число частиц фиксировано, при конечной температуре имеем
%
\begin{equation*}
    \varepsilon_F^{3/2}
    =
    \mu^{3/2}
    \left[
        1
        +
        \frac{1}{8}
        \left(
            \frac{\pi T}{\mu}
        \right)^2
    \right].
\end{equation*}

Разделим обе части на $\varepsilon_F^{3/2}$:
%
\begin{equation*}
    1
    =
    \left(
        \frac{\mu}{\varepsilon_F}
    \right)^{3/2}
    \left[
        1
        +
        \frac{1}{8}
        \left(
            \frac{\pi T}{\mu}
        \right)^2
    \right].
\end{equation*}

При
%
\begin{equation*}
    \mu
    \simeq
    \varepsilon_F
    \gg
    T
\end{equation*}
%
представим химический потенциал в виде
%
\begin{equation*}
    \mu
    =
    \varepsilon_F
    (1+\delta),
    \qquad
    |\delta|
    \ll
    1.
\end{equation*}

С требуемой точностью в температурной поправке можно заменить
$\mu$ на $\varepsilon_F$:
%
\begin{equation*}
    1
    =
    (1+\delta)^{3/2}
    \left[
        1
        +
        \frac{1}{8}
        \left(
            \frac{\pi T}{\varepsilon_F}
        \right)^2
    \right].
\end{equation*}

Разлагаем
%
\begin{equation*}
    (1+\delta)^{3/2}
    \simeq
    1
    +
    \frac{3}{2}\delta.
\end{equation*}

Следовательно,
%
\begin{equation*}
    1
    =
    1
    +
    \frac{3}{2}\delta
    +
    \frac{1}{8}
    \left(
        \frac{\pi T}{\varepsilon_F}
    \right)^2.
\end{equation*}

Отсюда
%
\begin{equation*}
    \frac{3}{2}\delta
    =
    -
    \frac{1}{8}
    \left(
        \frac{\pi T}{\varepsilon_F}
    \right)^2,
\end{equation*}
%
то есть
%
\begin{equation*}
    \delta
    =
    -
    \frac{1}{12}
    \left(
        \frac{\pi T}{\varepsilon_F}
    \right)^2.
\end{equation*}

Таким образом,
%
\[
    \rboxed{
        \mu
        =
        \varepsilon_F
        \left[
            1
            -
            \frac{1}{12}
            \left(
                \frac{\pi T}{\varepsilon_F}
            \right)^2
        \right]
    }.
\]

При увеличении температуры химический потенциал вырожденного
ферми-газа уменьшается.


% ------------------------------------------------------------
\subsection{Энтропия и теплоёмкость}
% ------------------------------------------------------------

Энтропия определяется выражением
%
\begin{equation*}
    S
    =
    -
    \left(
        \frac{\partial\Omega}{\partial T}
    \right)_{\mu,V}.
\end{equation*}

Из выражения для $\Omega$ получаем
%
\begin{equation*}
    S
    =
    \frac{\pi^2}{3}
    AV\sqrt{\mu}\,T.
\end{equation*}

В главном приближении по температуре
%
\begin{equation*}
    \mu
    \simeq
    \varepsilon_F,
\end{equation*}
%
поэтому
%
\begin{equation*}
    S
    =
    \frac{\pi^2}{3}
    AV\sqrt{\varepsilon_F}\,T.
\end{equation*}

Из соотношения
%
\begin{equation*}
    N
    =
    \frac{2}{3}
    AV\varepsilon_F^{3/2}
\end{equation*}
%
следует
%
\begin{equation*}
    AV\sqrt{\varepsilon_F}
    =
    \frac{3}{2}
    \frac{N}{\varepsilon_F}.
\end{equation*}

Следовательно,
%
\[
    \rboxed{
        S
        =
        N
        \frac{\pi^2}{2}
        \frac{T}{\varepsilon_F}
    }.
\]

Теплоёмкость при постоянном объёме
%
\begin{equation*}
    C_V
    =
    T
    \left(
        \frac{\partial S}{\partial T}
    \right)_V.
\end{equation*}

Поскольку при низких температурах
%
\begin{equation*}
    S
    \propto
    T,
\end{equation*}
%
получаем
%
\[
    \rboxed{
        C_V
        =
        \frac{\pi^2}{2}
        N
        \frac{T}{\varepsilon_F}
    }.
\]

Таким образом, теплоёмкость сильно вырожденного ферми-газа линейна
по температуре:
%
\begin{equation*}
    C_V
    \sim
    T.
\end{equation*}


% ------------------------------------------------------------
\subsection{Давление}
% ------------------------------------------------------------

Давление связано с $\Omega$-потенциалом соотношением
%
\begin{equation*}
    p
    =
    -
    \frac{\Omega}{V}.
\end{equation*}

Поэтому
%
\begin{equation*}
    p
    =
    \frac{2}{3}A
    \frac{2}{5}
    \mu^{5/2}
    \left[
        1
        +
        \frac{5}{8}
        \left(
            \frac{\pi T}{\mu}
        \right)^2
    \right].
\end{equation*}

Подставим
%
\begin{equation*}
    \mu
    =
    \varepsilon_F
    \left[
        1
        -
        \frac{1}{12}
        \left(
            \frac{\pi T}{\varepsilon_F}
        \right)^2
    \right].
\end{equation*}

Тогда
%
\begin{align*}
    p
    &=
    \frac{2}{3}A
    \frac{2}{5}
    \varepsilon_F^{5/2}
    \left[
        1
        -
        \frac{1}{12}
        \left(
            \frac{\pi T}{\varepsilon_F}
        \right)^2
    \right]^{5/2}
    \left[
        1
        +
        \frac{5}{8}
        \left(
            \frac{\pi T}{\varepsilon_F}
        \right)^2
    \right].
\end{align*}

Разложим первый множитель:
%
\begin{equation*}
    \left[
        1
        -
        \frac{1}{12}
        \left(
            \frac{\pi T}{\varepsilon_F}
        \right)^2
    \right]^{5/2}
    \simeq
    1
    -
    \frac{5}{24}
    \left(
        \frac{\pi T}{\varepsilon_F}
    \right)^2.
\end{equation*}

Поэтому
%
\begin{align*}
    p
    &=
    \frac{2}{3}A
    \frac{2}{5}
    \varepsilon_F^{5/2}
    \left[
        1
        -
        \frac{5}{24}
        \left(
            \frac{\pi T}{\varepsilon_F}
        \right)^2
        +
        \frac{5}{8}
        \left(
            \frac{\pi T}{\varepsilon_F}
        \right)^2
    \right]
    \\[6pt]
    &=
    \frac{2}{3}A
    \frac{2}{5}
    \varepsilon_F^{5/2}
    \left[
        1
        +
        \frac{5}{12}
        \left(
            \frac{\pi T}{\varepsilon_F}
        \right)^2
    \right].
\end{align*}

Так как
%
\begin{equation*}
    \frac{N}{V}
    =
    \frac{2}{3}A
    \varepsilon_F^{3/2},
\end{equation*}
%
окончательно
%
\[
    \rboxed{
        p
        =
        \frac{N}{V}
        \frac{2}{5}
        \varepsilon_F
        \left[
            1
            +
            \frac{5}{12}
            \left(
                \frac{\pi T}{\varepsilon_F}
            \right)^2
        \right]
    }.
\]

При $T=0$ получаем уже найденное ранее давление вырождения
%
\begin{equation*}
    p_0
    =
    \frac{2}{5}
    \frac{N}{V}
    \varepsilon_F.
\end{equation*}


% ------------------------------------------------------------
\subsection{Энергия}
% ------------------------------------------------------------

Для нерелятивистского идеального газа выполняется
%
\begin{equation*}
    E
    =
    \frac{3}{2}pV.
\end{equation*}

Следовательно,
%
\[
    \rboxed{
        E
        =
        \frac{3}{5}
        N\varepsilon_F
        \left[
            1
            +
            \frac{5}{12}
            \left(
                \frac{\pi T}{\varepsilon_F}
            \right)^2
        \right]
    }.
\]

При $T=0$
%
\begin{equation*}
    E_0
    =
    \frac{3}{5}
    N\varepsilon_F.
\end{equation*}


% ------------------------------------------------------------
\subsection{Зависимость давления от объёма}
% ------------------------------------------------------------

При фиксированном числе частиц
%
\begin{equation*}
    \varepsilon_F
    \sim
    \left(
        \frac{N}{V}
    \right)^{2/3}
    \sim
    V^{-2/3}.
\end{equation*}

Поэтому нулевой по температуре вклад в давление ведёт себя как
%
\begin{equation*}
    p_0
    \sim
    \frac{N}{V}\varepsilon_F
    \sim
    V^{-5/3}.
\end{equation*}

Температурная поправка имеет вид
%
\begin{equation*}
    \Delta p
    \sim
    p_0
    \frac{T^2}{\varepsilon_F^2}.
\end{equation*}

Так как
%
\begin{equation*}
    \varepsilon_F
    \sim
    V^{-2/3},
\end{equation*}
%
получаем
%
\begin{equation*}
    \Delta p
    \sim
    T^2V^{-1/3}.
\end{equation*}

Таким образом, при фиксированном $N$
%
\[
    \rboxed{
        p
        \sim
        V^{-5/3}
        +
        T^2V^{-1/3}
    }.
\]


% ============================================================
\section{Магнетизм электронного газа}
% ============================================================

Рассмотрим идеальный электронный газ во внешнем магнитном поле.

Для электрона спиновый $g$-фактор принимаем равным
%
\begin{equation*}
    g
    =
    2.
\end{equation*}

Введём магнетон Бора
%
\[
    \rboxed{
        \beta
        =
        \frac{
            e\hbar
        }{
            2m_ec
        }
    }.
\]

Во внешнем магнитном поле $H$ энергия электрона зависит от ориентации
его спина:
%
\[
    \rboxed{
        \varepsilon
        =
        \frac{p^2}{2m}
        \pm
        \beta H
    }.
\]

Таким образом, магнитное поле расщепляет энергетический спектр
на две ветви, соответствующие двум возможным ориентациям спина.


% ------------------------------------------------------------
\subsection{\texorpdfstring{$\Omega$}{Omega}-потенциал в магнитном поле}
% ------------------------------------------------------------

Для магнитной системы при фиксированном химическом потенциале
термодинамический дифференциал можно записать в виде
%
\begin{equation*}
    d\Omega
    =
    -S\,dT
    -
    p\,dV
    -
    \vec M\cdot d\vec H.
\end{equation*}

Обозначим через
%
\begin{equation*}
    \Omega_0
    =
    \left.
        \Omega
    \right|_{H=0}
\end{equation*}
%
термодинамический потенциал электронного газа в отсутствие
магнитного поля.

Число заполнения одночастичного состояния имеет вид
%
\begin{equation*}
    N_i
    =
    \left[
        \exp
        \left(
            \frac{
                p^2/(2m)
                \pm
                \beta H
                -
                \mu
            }{
                T
            }
        \right)
        +1
    \right]^{-1}.
\end{equation*}

Сдвиг энергии на $\pm\beta H$ эквивалентен противоположному
сдвигу химического потенциала для двух спиновых подсистем.

Поэтому
%
\begin{equation*}
    \Omega
    =
    \frac{1}{2}
    \Omega_0(\mu-\beta H)
    +
    \frac{1}{2}
    \Omega_0(\mu+\beta H).
\end{equation*}

Рассматривая слабое поле, разложим оба слагаемых по степеням $H$:
%
\begin{align*}
    \Omega_0(\mu-\beta H)
    &=
    \Omega_0(\mu)
    -
    \frac{\partial\Omega_0}{\partial\mu}
    \beta H
    +
    \frac{1}{2}
    \frac{\partial^2\Omega_0}{\partial\mu^2}
    \beta^2H^2
    +
    \ldots,
    \\[6pt]
    \Omega_0(\mu+\beta H)
    &=
    \Omega_0(\mu)
    +
    \frac{\partial\Omega_0}{\partial\mu}
    \beta H
    +
    \frac{1}{2}
    \frac{\partial^2\Omega_0}{\partial\mu^2}
    \beta^2H^2
    +
    \ldots.
\end{align*}

После сложения линейные по полю члены сокращаются:
%
\begin{equation*}
    \Omega
    =
    \Omega_0
    +
    \frac{1}{2}
    \frac{\partial^2\Omega_0}{\partial\mu^2}
    \beta^2H^2.
\end{equation*}

Но
%
\begin{equation*}
    N
    =
    -
    \frac{\partial\Omega_0}{\partial\mu},
\end{equation*}
%
следовательно,
%
\begin{equation*}
    \frac{\partial^2\Omega_0}{\partial\mu^2}
    =
    -
    \frac{\partial N}{\partial\mu}.
\end{equation*}

Таким образом,
%
\[
    \rboxed{
        \Omega
        =
        \Omega_0
        -
        \frac{1}{2}
        \left(
            \frac{\partial N}{\partial\mu}
        \right)_{T,V}
        \beta^2H^2
    }.
\]


% ------------------------------------------------------------
\subsection{Магнитный момент и восприимчивость}
% ------------------------------------------------------------

Магнитный момент определяется как
%
\begin{equation*}
    M
    =
    -
    \left(
        \frac{\partial\Omega}{\partial H}
    \right)_{T,V,\mu}.
\end{equation*}

Поэтому
%
\[
    \rboxed{
        M
        =
        \left(
            \frac{\partial N}{\partial\mu}
        \right)_{T,V}
        \beta^2H
    }.
\]

Магнитная восприимчивость равна
%
\begin{equation*}
    \chi
    =
    \frac{\partial M}{\partial H},
\end{equation*}
%
откуда
%
\[
    \rboxed{
        \chi
        =
        \left(
            \frac{\partial N}{\partial\mu}
        \right)_{T,V}
        \beta^2
    }.
\]

Найдём производную $\partial N/\partial\mu$ при
%
\begin{equation*}
    T
    \ll
    \varepsilon_F.
\end{equation*}

Ранее было получено
%
\begin{equation*}
    N
    =
    \frac{2}{3}AV
    \left[
        \mu^{3/2}
        +
        \frac{(\pi T)^2}{8\sqrt{\mu}}
    \right].
\end{equation*}

Дифференцируя по $\mu$, имеем
%
\begin{align*}
    \frac{\partial N}{\partial\mu}
    &=
    \frac{2}{3}AV
    \left[
        \frac{3}{2}\mu^{1/2}
        -
        \frac{1}{16}
        \frac{(\pi T)^2}{\mu^{3/2}}
    \right]
    \\[6pt]
    &=
    AV\sqrt{\mu}
    \left[
        1
        -
        \frac{1}{24}
        \left(
            \frac{\pi T}{\mu}
        \right)^2
    \right].
\end{align*}

При этом
%
\begin{equation*}
    \mu
    =
    \varepsilon_F
    \left[
        1
        -
        \frac{1}{12}
        \left(
            \frac{\pi T}{\varepsilon_F}
        \right)^2
    \right].
\end{equation*}

Следовательно,
%
\begin{align*}
    \sqrt{\mu}
    &=
    \sqrt{\varepsilon_F}
    \left[
        1
        -
        \frac{1}{12}
        \left(
            \frac{\pi T}{\varepsilon_F}
        \right)^2
    \right]^{1/2}
    \\[4pt]
    &\simeq
    \sqrt{\varepsilon_F}
    \left[
        1
        -
        \frac{1}{24}
        \left(
            \frac{\pi T}{\varepsilon_F}
        \right)^2
    \right].
\end{align*}

Поэтому
%
\begin{align*}
    \frac{\partial N}{\partial\mu}
    &=
    AV\sqrt{\varepsilon_F}
    \left[
        1
        -
        \frac{1}{24}
        \left(
            \frac{\pi T}{\varepsilon_F}
        \right)^2
    \right]
    \left[
        1
        -
        \frac{1}{24}
        \left(
            \frac{\pi T}{\varepsilon_F}
        \right)^2
    \right]
    \\[6pt]
    &=
    AV\sqrt{\varepsilon_F}
    \left[
        1
        -
        \frac{1}{12}
        \left(
            \frac{\pi T}{\varepsilon_F}
        \right)^2
    \right].
\end{align*}

Используем
%
\begin{equation*}
    N
    =
    \frac{2}{3}
    AV\varepsilon_F^{3/2}.
\end{equation*}

Отсюда
%
\begin{equation*}
    AV\sqrt{\varepsilon_F}
    =
    \frac{3}{2}
    \frac{N}{\varepsilon_F}.
\end{equation*}

Таким образом,
%
\[
    \rboxed{
        \frac{\partial N}{\partial\mu}
        =
        \frac{3}{2}
        \frac{N}{\varepsilon_F}
        \left[
            1
            -
            \frac{1}{12}
            \left(
                \frac{\pi T}{\varepsilon_F}
            \right)^2
        \right]
    }.
\]

Следовательно, магнитный момент электронного газа при
$T\ll\varepsilon_F$ равен
%
\[
    \rboxed{
        M
        =
        \frac{3}{2}
        \frac{N}{\varepsilon_F}
        \left[
            1
            -
            \frac{1}{12}
            \left(
                \frac{\pi T}{\varepsilon_F}
            \right)^2
        \right]
        \beta^2H
    }.
\]

Магнитная восприимчивость имеет вид
%
\[
    \rboxed{
        \chi
        =
        \frac{3}{2}
        \frac{N}{\varepsilon_F}
        \left[
            1
            -
            \frac{1}{12}
            \left(
                \frac{\pi T}{\varepsilon_F}
            \right)^2
        \right]
        \beta^2
    }.
\]

Так как
%
\begin{equation*}
    \chi
    >
    0,
\end{equation*}
%
электронный газ обладает \rconcept{парамагнетизмом Паули}.

При абсолютном нуле
%
\[
    \rboxed{
        \chi(0)
        =
        \frac{
            3N\beta^2
        }{
            2\varepsilon_F
        }
    }.
\]

При повышении температуры в области $T\ll\varepsilon_F$
парамагнитная восприимчивость уменьшается как
%
\begin{equation*}
    \chi(T)
    =
    \chi(0)
    \left[
        1
        -
        \frac{1}{12}
        \left(
            \frac{\pi T}{\varepsilon_F}
        \right)^2
    \right].
\end{equation*}

В противоположном пределе
%
\begin{equation*}
    T
    \gg
    \varepsilon_F
\end{equation*}
%
электронный газ переходит в невырожденный, больцмановский режим.

\end{document}
